% GENERATED FILE. Edit canonical lessons, then run npm run generate:books.
% canonical-source-hash: fnv1a64:4e476fc99fdaa6c9
% canonical-lessons: TA-C165-palaku, TA-C165-kontatu, TA-C165-alankari, TA-C165-natanamatu, TA-C165-vananku

\chapter{To get to know, To celebrate, To decorate, To dance, To worship}
\label{ch:ta-to-get-to-know-to-celebrate-to-decorate-to-dance-to-worship}
\hlchaptermodality{165}

\begin{chapteropening}
\textbf{By the end of this chapter:} \emph{I can say to get to know, to celebrate, to decorate, to dance and to worship.}

Complete the last lesson of chapter 165: I can say to get to know, to celebrate, to decorate, to dance and to worship.
\end{chapteropening}

\section[\texorpdfstring{\ta{பழகு} (paḻaku)}{பழகு (paḻaku)}]{\ta{பழகு} (paḻaku) — to get to know, to mix with (people)}

\label{lesson:TA-C165-palaku}



\begin{quote}
Before the new one: say the Tamil for to pass, then the Tamil for to pay attention.
\end{quote}

\subsection*{You'll want to know: \ta{பழகு}}
\textbf{\ta{பழகு}} — \emph{paḻaku} — \textquotedblleft{}to get to know, to mix with (people)\textquotedblright{}.

\begin{grammarlens}[title={What you've built}]
One more word for this chapter.
\end{grammarlens}

\subsection*{Your turn}
\begin{itemize}
\raggedright
  \item \emph{Say it:} \emph{paḻaku}
  \item \emph{Say it:} \emph{paḻaku}, once more
  \item \emph{Say it:} say \emph{paḻaku}
  \item \emph{Recall:} say the Tamil for to pass, then the Tamil for to pay attention, then say \emph{paḻaku} again
\end{itemize}

\begin{tcolorbox}[breakable,colback=teal!4,colframe=teal!35!black,title={Before you move on}]
What does \ta{பழகு} mean? (\textquotedblleft{}To get to know, to mix with (people)\textquotedblright{}.) Say it once more.
\end{tcolorbox}
\section[\texorpdfstring{\ta{கொண்டாடு} (koṇṭāṭu)}{கொண்டாடு (koṇṭāṭu)}]{\ta{கொண்டாடு} (koṇṭāṭu) — to celebrate}

\label{lesson:TA-C165-kontatu}



\begin{quote}
Before the new one: say the Tamil for to pay attention, then the Tamil for to get to know.
\end{quote}

\subsection*{You'll want to know: \ta{கொண்டாடு}}
\textbf{\ta{கொண்டாடு}} — \emph{koṇṭāṭu} — \textquotedblleft{}to celebrate\textquotedblright{}.

\begin{grammarlens}[title={What you've built}]
One more word for this chapter.
\end{grammarlens}

\subsection*{Your turn}
\begin{itemize}
\raggedright
  \item \emph{Say it:} \emph{koṇṭāṭu}
  \item \emph{Say it:} \emph{koṇṭāṭu}, once more
  \item \emph{Say it:} say \emph{koṇṭāṭu}
  \item \emph{Recall:} say the Tamil for to pay attention, then the Tamil for to get to know, then say \emph{koṇṭāṭu} again
\end{itemize}

\begin{tcolorbox}[breakable,colback=teal!4,colframe=teal!35!black,title={Before you move on}]
What does \ta{கொண்டாடு} mean? (\textquotedblleft{}To celebrate\textquotedblright{}.) Say it once more.
\end{tcolorbox}
\section[\texorpdfstring{\ta{அலங்கரி} (alaṅkari)}{அலங்கரி (alaṅkari)}]{\ta{அலங்கரி} (alaṅkari) — to decorate}

\label{lesson:TA-C165-alankari}



\begin{quote}
Before the new one: say the Tamil for to get to know, then the Tamil for to celebrate.
\end{quote}

\subsection*{You'll want to know: \ta{அலங்கரி}}
\textbf{\ta{அலங்கரி}} — \emph{alaṅkari} — \textquotedblleft{}to decorate\textquotedblright{}.

\begin{grammarlens}[title={What you've built}]
One more word for this chapter.
\end{grammarlens}

\subsection*{Your turn}
\begin{itemize}
\raggedright
  \item \emph{Say it:} \emph{alaṅkari}
  \item \emph{Say it:} \emph{alaṅkari}, once more
  \item \emph{Say it:} say \emph{alaṅkari}
  \item \emph{Recall:} say the Tamil for to get to know, then the Tamil for to celebrate, then say \emph{alaṅkari} again
\end{itemize}

\begin{tcolorbox}[breakable,colback=teal!4,colframe=teal!35!black,title={Before you move on}]
What does \ta{அலங்கரி} mean? (\textquotedblleft{}To decorate\textquotedblright{}.) Say it once more.
\end{tcolorbox}
\section[\texorpdfstring{\ta{நடனமாடு} (naṭaṉamāṭu)}{நடனமாடு (naṭaṉamāṭu)}]{\ta{நடனமாடு} (naṭaṉamāṭu) — to dance}

\label{lesson:TA-C165-natanamatu}



\begin{quote}
Before the new one: say the Tamil for to celebrate, then the Tamil for to decorate.
\end{quote}

\subsection*{You'll want to know: \ta{நடனமாடு}}
\textbf{\ta{நடனமாடு}} — \emph{naṭaṉamāṭu} — \textquotedblleft{}to dance\textquotedblright{}.

\begin{grammarlens}[title={What you've built}]
One more word for this chapter.
\end{grammarlens}

\subsection*{Your turn}
\begin{itemize}
\raggedright
  \item \emph{Say it:} \emph{naṭaṉamāṭu}
  \item \emph{Say it:} \emph{naṭaṉamāṭu}, once more
  \item \emph{Say it:} say \emph{naṭaṉamāṭu}
  \item \emph{Recall:} say the Tamil for to celebrate, then the Tamil for to decorate, then say \emph{naṭaṉamāṭu} again
\end{itemize}

\begin{tcolorbox}[breakable,colback=teal!4,colframe=teal!35!black,title={Before you move on}]
What does \ta{நடனமாடு} mean? (\textquotedblleft{}To dance\textquotedblright{}.) Say it once more.
\end{tcolorbox}
\section[\texorpdfstring{\ta{வணங்கு} (vaṇaṅku)}{வணங்கு (vaṇaṅku)}]{\ta{வணங்கு} (vaṇaṅku) — to worship, to bow to}

\label{lesson:TA-C165-vananku}



\begin{quote}
Before the new one: say the Tamil for to decorate, then the Tamil for to dance.
\end{quote}

\subsection*{You'll want to know: \ta{வணங்கு}}
\textbf{\ta{வணங்கு}} — \emph{vaṇaṅku} — \textquotedblleft{}to worship, to bow to\textquotedblright{}.

\begin{grammarlens}[title={What you've built}]
One more word for this chapter.
\end{grammarlens}

\subsection*{Your turn}
\begin{itemize}
\raggedright
  \item \emph{Say it:} \emph{vaṇaṅku}
  \item \emph{Say it:} \emph{vaṇaṅku}, once more
  \item \emph{Say it:} say \emph{vaṇaṅku}
  \item \emph{Recall:} say the Tamil for to decorate, then the Tamil for to dance, then say \emph{vaṇaṅku} again
\end{itemize}

\begin{tcolorbox}[breakable,colback=teal!4,colframe=teal!35!black,title={Before you move on}]
What does \ta{வணங்கு} mean? (\textquotedblleft{}To worship, to bow to\textquotedblright{}.) Say it once more.
\end{tcolorbox}
